\documentclass[10pt,a4paper]{amsart}
\usepackage[margin=19mm]{geometry}
\usepackage{amsmath,amssymb,mathtools}
\usepackage[scr=rsfs,cal=euler]{mathalfa}
\usepackage{xfrac}
\usepackage[unicode,hidelinks,bookmarks=false]{hyperref}
\newcommand{\ie}{i.e.\ }
\newcommand{\cf}{cf.\ }
\newcommand{\resp}{respectively\ }
\newcommand{\Subsection}{\subsection}
\providecommand{\nobreakdash}{\nobreak\hbox{-}}
\setlength{\emergencystretch}{2em}
\multlinegap0pt
\usepackage{amssymb}
\usepackage{xcolor}
\usepackage{tikz}
\usepackage[normalem]{ulem}
\usepackage{dsfont}
\usepackage[shortlabels]{enumitem}
\newtheorem{proposition}{Proposition}
\newcommand{\Deltalb}{\Delta_{\IS^2}}
\newcommand{\E}{{\mathbb E}}
\newcommand{\IS}{\mathbb{S}}
\newcommand{\R}{{\mathbb R}}
\DeclareMathOperator{\Var}{Var}
\newcommand{\cE}{\mathcal E}
\newcommand{\cP}{\mathcal{P}}
\newcommand{\tr}{{\operatorname{Tr}\, }}
\title{Long-time behavior of exact and numerical solutions of stochastic evolution equations on the sphere}
\author{David Cohen, Bj\"orn M\"uller, Andrea Papini}
\date{Source: arXiv:2604.05644v1 (CC BY 4.0)}
\begin{document}
\maketitle

\noindent\emph{Source and licence.} Long-time behavior of exact and numerical solutions of stochastic evolution equations on the sphere. Version arXiv:2604.05644v1 (CC BY 4.0), distributed by arXiv under the Creative Commons Attribution 4.0 International licence, http://creativecommons.org/licenses/by/4.0/, as recorded in the arXiv licence metadata for this exact version and shown on the abstract page https://arxiv.org/abs/2604.05644v1. Full licence notice: \texttt{licenses/arxiv-2604.05644-CC-BY-4.0.txt}.\par\smallskip

\noindent\textit{Editorial scope.} Counted unit: the article's proposition prop:energy-bem (source0.tex lines 718--732). The article's complete original proof of it is delivered with it (source0.tex lines 733--764, one \texttt{proof} environment of 32 lines). No mathematics is written, repaired or re-derived: the statement and the proof are the article's own lines, in the article's own notation, and no retained line is substituted or re-flowed.  Retained uncounted as the source-local context the proof needs: lines 371--373; lines 462--464; lines 468--470; lines 596--603; lines 634--636; lines 643--653; lines 697--715.  Nothing was omitted from any retained span, so the package carries no omission record.  No retained line contains a citation, so the wrapper carries no bibliography and no bibliography entry.  No retained line contains a figure environment, an image-inclusion command or drawing code, so no image is delivered and the package declares no asset dependency.  Editorial formatting only, in addition: an AMS \texttt{amsart} preamble that restates the article's own declarations and macros used by the retained lines, namely the environments \texttt{proposition} on separate counters, together with the standard packages those lines need.  The retained environments print in this document as Proposition 1 in order of appearance, whereas the article prints the same items with the numbering of its own full document.  The complete native source of the article as distributed in the arXiv e-print for this exact version is bundled unchanged as \texttt{sources/197970/source0.tex}.
\begin{equation}\label{swe}
\partial_{tt}u(t)-\Deltalb u(t)=\dot{L}(t)
\end{equation}

\begin{equation}\label{eEM}
X^{\text{EM}}_{n+1}=X^{\text{EM}}_{n}+\tau A^\kappa X^{\text{EM}}_{n}+\cP_\kappa\left(B\Delta L_n\right).
\end{equation}

\begin{equation}\label{bEM}
X^{\text{bEM}}_{n+1}=X^{\text{bEM}}_{n}+\tau A^\kappa X^{\text{bEM}}_{n+1}+\cP_\kappa\left(B\Delta L_n\right).
\end{equation}

We start by investigating the long-time behavior, with respect to the energy, of the forward Euler--Maruyama scheme~\eqref{eEM}. When applied to the stochastic wave equation on the sphere~\eqref{swe}, this time integrator reads 
\begin{align}
    \begin{split}
\label{eqn:SWE_EM}
u_{1,n+1}^\kappa&=u_{1,n}^\kappa+\tau u_{2,n}^\kappa,\\
u_{2,n+1}^\kappa&=u_{2,n}^\kappa+\tau \Deltalb^\kappa u_{1,n}^\kappa+\cP_\kappa\Delta L_n.
    \end{split}
\end{align}

\begin{equation}\label{energEM}
\mathcal{E}^\kappa_n=\mathcal{E}^{0,0}_n+\sum_{\ell=1}^{\kappa}\sum_{m=-\ell}^{\ell} \mathcal{E}^{\ell,m}_n,
\end{equation}

\begin{proposition}\label{prop:energy-em}
Consider the stochastic wave equation on the sphere~\eqref{swe} with initial values $(u(0),\partial_tu(0))=(v_1,v_2)\in [L^2(\IS^2,\R)]^2$.
Assume that the expected value of the initial energy $\E\left[\mathcal{E}(0)\right]<\infty$ and that the $Q$-L\'evy process $L$ is trace-class.
Furthermore, assume that the expected value of the initial energy satisfies $0<\E\left[\mathcal{E}^\kappa_0-\cE_0^{0,0}\right]$.
Let $(u_{1,n}^\kappa, u_{2,n}^\kappa)$ denote the fully-discrete approximation given by the forward Euler--Maruyama scheme~\eqref{eqn:SWE_EM}.
The energy of this fully-discrete approximation grows exponentially fast in time:
\begin{equation*}
\E\left[\mathcal{E}^\kappa_n\right] \geq \exp\left(\frac12\tau t_n\right)\E\left[\mathcal{E}^\kappa_0 - \mathcal{E}_0^{0,0}\right]
\end{equation*}
for every discrete times $t_n=n\tau$ with integers $n\geq0$.
\end{proposition}

\subsection{Backward Euler--Maruyama method}
Our next goal is to investigate the long-time behavior, with respect to the energy, of the backward Euler--Maruyama scheme~\eqref{bEM} when applied to the stochastic wave equation on the sphere~\eqref{swe}. This numerical scheme is given by the following implicit relation
\begin{align}
\begin{split}
    \label{eqn:SWE_BEM}
    u_{1,n+1}^\kappa&=u_{1,n}^\kappa+\tau u_{2,n+1}^\kappa,\\
    u_{2,n+1}^\kappa&=u_{2,n}^\kappa+\tau \Deltalb^\kappa u_{1,n+1}^\kappa+\cP_\kappa\Delta L_n.
\end{split}
\end{align}
As we did previously for the forward Euler--Maruyama scheme, we obtain the following system of equation for the backward Euler--Maruyama scheme~\eqref{eqn:SWE_BEM}
\begin{align}
    \begin{split}
    \label{eqn:SWE_BEM2}
    u_{1,n+1}^{\ell, m} &= u_{1, n}^{\ell, m} + \tau u_{2, n+1}^{\ell, m}\\
    \left(1+\tau^2\lambda_\ell\right)u_{2,n+1}^{\ell, m} &= u_{2, n}^{\ell, m} - \tau\lambda_\ell u_{1, n}^{\ell, m} + \Delta L^{\ell, m}_n
    \end{split}
\end{align}
for $\ell=0,\ldots,\kappa$, $m=-\ell,\ldots,\ell$, and $n\geq 0$. We use the same notation for the energy of the backward Euler--Maruyama scheme as
for the forward Euler--Maruyama scheme, see~\eqref{energEM}.

\begin{proposition}\label{prop:energy-bem}
Consider the stochastic wave equation on the sphere~\eqref{swe} with initial values $(u(0),\partial_tu(0))=(v_1,v_2)\in [L^2(\IS^2,\R)]^2$.
Assume that the expected value of the initial energy $\E\left[\mathcal{E}(0)\right]<\infty$ and that the $Q$-L\'evy process $L$ is trace-class.
Furthermore, assume that the expected value of the initial energy satisfies $0<\E\left[\mathcal{E}^\kappa_0\right]$.
Let $\left(u_{1,n}^\kappa, u_{2,n}^\kappa\right)$ denote the fully-discrete approximation given by the backward Euler--Maruyama scheme~\eqref{eqn:SWE_BEM}.
The energy of this fully-discrete approximation grows at a slower rate than the exact solution to the considered SPDE:
\begin{equation*}
		\E\left[\cE^\kappa_n\right] <
         \E\left[\cE^\kappa_0\right]+ \frac{1}{2\tau}\tr\left(Q^\kappa\right) + \frac {t_n} {2} v_{0,0}
\end{equation*}
for every discrete times $t_n=n\tau$ with integers $n\geq0$ and where we recall that $v_{0,0}=\Var(L_{0,0}(1))$. More so, we have
\begin{equation}\label{eqn:SWE_energy_BEM_limit}
    \lim_{t_n\to\infty}\frac{\E\left[\mathcal{E}^\kappa_n\right]}{t_n}=\frac 1 2 v_{0,0}.
\end{equation}
\end{proposition}

\begin{proof}	
	Let $\ell\geq 1$. As in the proof of Proposition~\ref{prop:energy-em}, we substitute \eqref{eqn:SWE_BEM2} in the $(\ell, m)$ mode of the energy and obtain:
	\begin{align*}
		\E\left[\cE^{\ell, m}_{n+1}\right] &= \frac{1}{2}\E\left[\left(u^{\ell, m}_{2, n+1}\right)^2 + \lambda_\ell \left(u^{\ell, m}_{1, n+1}\right)^2\right]=\left(1 + \tau^2\lambda_\ell\right)^{-1} \left(\E\left[\cE^{\ell, m}_{n}\right] + \frac{1}2v_{\ell, m}\tau\right).
	\end{align*}
	We again observe that for the mode $\ell = 0$, the energy is preserved exactly.
	However, for all other modes, energy is lost for the backward Euler--Maruyama scheme.
	
	In particular, we have that
	\[
	\E\left[\cE^{\ell, m}_{n+1}\right] \leq \left(1 + \tau^2\right)^{-1} \left(\E\left[\cE^{\ell, m}_{n}\right] + \frac{1}2v_{\ell, m}\tau\right).
	\]
	Iterating this for the $n$--th increment implies
	\begin{align*}
		\E\left[\cE^{\ell, m}_{n}\right] &\leq \left(1+\tau^2\right)^{-n} \E\left[\cE^{\ell, m}_{0}\right] + \sum_{k=1}^{n}\left(1+\tau^2\right)^{-k}\frac{1}{2}v_{\ell, m}\tau\\
		& \leq (1+\tau^2)^{-n} \E\left[\cE^{\ell, m}_{0}\right]  + \frac{\tau}{2}v_{\ell,m} \left(\frac{1-\left(1 + \tau^2\right)^{-n} }{\tau^2}\right)\\
		& \leq (1+\tau^2)^{-n} \E\left[\cE^{\ell, m}_{0}\right]  + \frac{1}{2\tau}v_{\ell,m}.
	\end{align*}
	
	Summing across all modes, this gives directly the evolution of the expected energy:
	\begin{align}\label{bem:bound}
		\E\left[\cE^\kappa_n\right] &\leq \left(1+\tau^2\right)^{-n} \E\left[\cE^\kappa_0 - \cE_0^{0, 0}\right]+ \frac{1}{2\tau}\tr\left(Q^\kappa\right) + \E\left[\cE_n^{0,0}\right],
	\end{align}
    where we note that $\tr\left(Q^\kappa\right)=\sum_{\ell=0}^{\kappa}\sum_{m=-\ell}^{\ell} v_{\ell,m}$.
	This clearly shows that the energy of the numerical solution grows more slowly than the energy of the true solution. 

    Moreover, one observes that the inequality \eqref{bem:bound} implies the estimates
	\begin{equation*}
		\E\left[\cE^\kappa_n\right] \leq (1+\tau^2)^{-n} \E\left[\cE^\kappa_0\right]+ \frac{1}{2\tau}\tr\left(Q^\kappa\right) + \frac {t_n} {2} v_{0,0} < \E\left[\cE^\kappa_0\right]+ \frac{1}{2\tau}\tr\left(Q^\kappa\right) + \frac {t_n} {2} v_{0,0}.
	\end{equation*}
	This concludes the proof of the proposition.
\end{proof}

\end{document}
